\documentclass[11pt]{article}
\usepackage[spanish]{babel}
\usepackage[utf8]{inputenc}
\usepackage[T1]{fontenc}
\usepackage{booktabs}
\usepackage{graphicx}
\usepackage{hyperref}
\usepackage[a4paper,margin=2.5cm]{geometry}

\title{Lab 01: Métodos no supervisados}
\author{Nombre completo (ID) \\ Nombre completo (ID)}
\date{EIF420O -- I Ciclo 2026}

\begin{document}
\maketitle

% Reemplace esta estructura por la plantilla oficial del profesor si se proporciona.
\section{Introducción y objetivos}
Describa el conjunto de datos provisto por el profesor, su procedencia, tamaño,
variables y el objetivo del análisis.

\section{Fundamentos teóricos}
El EDA permite caracterizar distribuciones, valores faltantes, atípicos y
relaciones entre variables antes de modelar. Esta etapa define decisiones de
limpieza, codificación y escala que deben conservarse para reproducir los
experimentos.

El ACP transforma variables posiblemente correlacionadas en componentes
ortogonales que retienen la mayor varianza posible \cite{jolliffe2002}. Para
cada configuración se reporta la varianza explicada y acumulada, junto con la
interpretación de las cargas.

K-Means asigna observaciones a $k$ centroides minimizando la suma de cuadrados
intra-grupo \cite{macqueen1967}. HAC construye una jerarquía de fusiones; el
enlace elegido determina cómo se calcula la distancia entre grupos
\cite{murtagh2012}. En ambos casos, silhouette es evidencia para comparar
configuraciones, no un sustituto de la interpretación de los grupos.

t-SNE conserva vecindades locales en una proyección de baja dimensionalidad
\cite{vandermaaten2008}; UMAP modela una estructura de vecindad para producir
una representación compacta \cite{mcinnes2018}. Sus gráficos son herramientas
exploratorias y no prueban por sí solos la existencia de clústeres.

\begin{center}
\fbox{\parbox{0.88\linewidth}{\centering
\textbf{Mapa conceptual:} datos $\rightarrow$ EDA $\rightarrow$
limpieza/codificación/escala $\rightarrow$
\{ACP, t-SNE, UMAP, K-Means, HAC\} $\rightarrow$
evaluación (varianza, silhouette, trustworthiness) $\rightarrow$
interpretación y selección del modelo.}}
\end{center}

\section{Metodología reproducible}
Indique versión de Python, dependencias, semilla, variables seleccionadas,
tratamiento de faltantes, codificación y estandarización. Adjunte el comando
utilizado, por ejemplo:
\begin{verbatim}
python -m scripts.ejecutar_lab01 data/datos_profesor.csv --salida salida_lab01 --semilla 42
\end{verbatim}

\section{Resultados y análisis}
\subsection{ACP}
Reporte las configuraciones de componentes, varianza explicada/acumulada y
analice la configuración seleccionada. Inserte sus figuras y tablas.

\subsection{K-Means}
Compare valores de $k$ con inercia y silhouette. Justifique el modelo elegido.

\subsection{HAC}
Compare enlaces y número de grupos; interprete el dendrograma y silhouette.

\subsection{t-SNE y UMAP}
Compare perplexity (t-SNE), vecinos y distancia mínima (UMAP). Use las
proyecciones y trustworthiness como evidencia de apoyo, no como única decisión.

\section{Conclusiones}
Sintetice los hallazgos, limitaciones y la justificación final de cada modelo.

\bibliographystyle{plain}
\bibliography{referencias}
\end{document}
